\documentclass[12pt]{article}
\usepackage[utf8]{inputenc}
\usepackage{amsmath,amsfonts,amssymb}
\usepackage{booktabs}
\usepackage{graphicx}
\usepackage{hyperref}
% set page margins and reduce top whitespace before title
\usepackage[margin=1in,top=0.8in]{geometry}
\title{HW3 --- Multi-Layer Perceptron: Answers and Results}
\author{Rahul Gonsalves}
\date{\today}

\begin{document}
\maketitle

\section*{Question 1: Loss functions (10 points)}
\subsection*{(a) Definitions and unified picture}
\textbf{0--1 loss.} For binary labels $y\in\{\pm1\}$ and score $f(x)$,
\[ L_{0-1}(y,f(x)) = \mathbf{1}[y f(x) \le 0]. \]
It counts misclassification; non-convex and not useful for gradient-based training.

\textbf{Hinge loss (SVM).}
\[ L_{\text{hinge}}(y,f(x)) = \max(0,\,1 - y f(x)). \]
Margin-based, convex, piecewise-linear. Loss is zero when margin $m=y f(x)\ge1$.

\textbf{Logistic loss.}
\[ L_{\text{log}}(y,f(x)) = \log(1 + e^{-y f(x)}). \]
Convex and smooth surrogate for 0--1; penalizes negative margins strongly, decays smoothly for positive margins.

\textbf{Unified picture.} All three are margin-based losses $\phi(m)$ with $m=yf(x)$. 0--1 is the ideal discrete loss. Hinge and logistic are convex surrogates that enable optimization; hinge becomes exactly zero for margin beyond a threshold, logistic only approaches zero asymptotically.

\subsection*{(b) Influence of a distant correctly-classified point}
If a point has large positive margin $m\gg1$: hinge loss has zero derivative (it is inactive) and the point does not affect the SVM solution (non-support vector). Logistic loss has derivative $-y/(1+e^{y f(x)})$, which is small but nonzero for finite $m$, so logistic regression is still (slightly) affected by such points. Intuition: SVM focuses on support vectors near the boundary; logistic regression ``weights'' all points.

\section*{Question 2: Backprop when output activation is identity (20 points)}
\subsection*{Formulas}
Notation (layer index $\ell$, output layer $L$):
\begin{align*}
s^{(\ell)} &= W^{(\ell)} x^{(\ell-1)} + b^{(\ell)},\\
x^{(\ell)} &= \phi^{(\ell)}(s^{(\ell)}),
\end{align*}
with the output activation identity: $x^{(L)} = s^{(L)}$. Loss for one example is
\[ e = \tfrac{1}{2} \|x^{(L)} - y\|^2. \]

Backprop results:
\begin{align*}
\delta^{(L)} &= \frac{\partial e}{\partial s^{(L)}} = x^{(L)} - y,\\
\delta^{(\ell)} &= (W^{(\ell+1)})^\top \delta^{(\ell+1)} \odot \phi'^{(\ell)}(s^{(\ell)}) \qquad (\ell < L),\\
\frac{\partial e}{\partial W^{(\ell)}} &= \delta^{(\ell)}\,(x^{(\ell-1)})^\top,\\
\frac{\partial e}{\partial b^{(\ell)}} &= \delta^{(\ell)}.
\end{align*}

\subsection*{Worked numeric example (one-hidden-layer, tanh)}
The numeric example from the assignment has been computed and is summarized below (values rounded):
\begin{itemize}
  \item $x^{(0)}=[1.0,2.0]^\top$, $y=0$.
  \item $W^{(1)}=\begin{bmatrix}0.1 & -0.2\\0.4 & 0.5\end{bmatrix}$, $b^{(1)}=[0.0,0.1]^\top$.
  \item $W^{(2)}=[0.3, -0.6]$, $b^{(2)}=0.05$.
\end{itemize}
Forward/backward results (rounded):
\begin{align*}
s^{(1)} &\approx [-0.300,\;1.500]^\top,\quad x^{(1)} \approx [-0.2913,\;0.9051]^\top,\\
s^{(2)} &= x^{(2)} \approx -0.58048,\\
\delta^{(2)} &\approx -0.58048,\\
\delta^{(1)} &\approx [-0.15937,\;0.06295]^\top,\\
\frac{\partial e}{\partial W^{(2)}} &\approx [0.16902,\;-0.52543],\quad \frac{\partial e}{\partial b^{(2)}}\approx -0.58048,\\
\frac{\partial e}{\partial W^{(1)}} &\approx \begin{bmatrix}-0.15937 & -0.31873\\[3pt]0.06295 & 0.12590\end{bmatrix},\quad
\frac{\partial e}{\partial b^{(1)}}\approx [-0.15937,\;0.06295]^\top.
\end{align*}

% A suggestion: if the grader prefers the full step-by-step arithmetic, include the per-scalar computations as shown in class.

\section*{Question 3: Parameter count (10 points)}
Network: input dimension $20$, five hidden layers each with $10$ units, scalar output.
\begin{itemize}
  \item Input$\to$H1: $W_1$ is $10\times20$ (200 weights) + 10 biases = 210.
  \item H1$\to$H2, H2$\to$H3, H3$\to$H4, H4$\to$H5: each $10\times10$ (100) + 10 biases = 110 per layer, four layers $\to$ 440.
  \item H5$\to$Output: $1\times10$ (10) + 1 bias = 11.
\end{itemize}
Total parameters $=210 + 440 + 11 = 661$.

\section*{Question 4: Programming part (60 points)}
\subsection*{Implementation summary}
I implemented the requested components in the provided starter code:
\begin{itemize}
  \item \texttt{code/DataReader.py}: implemented \texttt{prepare\_X} to produce three features per image:
    \begin{itemize}
      \item Feature 1 (symmetry): compute horizontal flip of the $16\times16$ image, take the mean absolute difference between original and flipped, divide by 256 and negate (so more symmetric images have larger/less-negative values).
      \item Feature 2 (intensity): mean pixel intensity computed as the sum over pixels divided by 256.
      \item Feature 3 (bias): constant 1.
    \end{itemize}
  \item \texttt{code/MLP.py}: implemented a 1-hidden-layer MLP using \texttt{nn.Linear} and ReLU (Linear $\to$ ReLU $\to$ Linear(3 logits)).
  \item \texttt{code/main.py}: implemented the \texttt{evaluation} function (loads the best validation checkpoint and computes test accuracy); made \texttt{HIDDEN\_SIZE} overridable via an environment variable.
\end{itemize}

\subsection*{Experimental setup}
Common hyperparameters used for all runs:
\begin{itemize}
  \item Epochs: 50; Batch size: 32; Optimizer: Adam (lr=0.001); Loss: CrossEntropyLoss.
  \item Seed: 42 (sets Python, NumPy, PyTorch seeds for reproducibility).
  \item Input features: 3 (bias, symmetry, intensity).
\end{itemize}

\subsection*{Hyperparameter tuning (hidden size)}
I trained and evaluated models with hidden sizes $H\in\{1,3,5,10,20\}$. Each run saved the best model according to validation accuracy and then reported test accuracy.

\begin{table}[h!]
\centering
\begin{tabular}{@{}cccc@{}}
\toprule
Hidden size $H$ & \# parameters & Best val acc & Test acc \\
\midrule
1  & $7\cdot1+3=10$  & 0.6968 & 0.6943 \\
3  & $7\cdot3+3=24$  & 0.6952 & 0.6906 \\
5  & $7\cdot5+3=38$  & 0.6619 & 0.8100 \\
10 & $7\cdot10+3=73$ & 0.7492 & 0.8551 \\
20 & $7\cdot20+3=143$& 0.7524 & 0.8575 \\
\bottomrule
\end{tabular}
\caption{Results for different hidden sizes. \# parameters computed by $7H+3$.}
\end{table}

\subsection*{Observations and analysis}
\begin{itemize}
  \item Small networks (H=1) underfit and achieve lower test accuracy. Even H=3 performs well, showing the three engineered features are informative.
  \item Increasing capacity to H=10 or H=20 yields modest improvements; best test accuracy observed was \textbf{0.8575} (H=20).
  \item Validation and test accuracies track closely; model selection by best validation accuracy produced good test performance.
  \item No strong overfitting was observed for these model sizes (validation accuracy improved alongside decreasing training loss).
\end{itemize}

\subsection*{Reproducibility and how to run}
Requirements: Python 3, \texttt{numpy}, \texttt{torch} (CPU wheel is sufficient).
\begin{verbatim}
# from repo root
cd code
python3 main.py                # run with default HIDDEN_SIZE=3
HIDDEN_SIZE=10 python3 main.py # change hidden size
\end{verbatim}
Best model is saved as \texttt{best\_mlp\_model.pth} (this file is ignored by .gitignore so it won't be committed).

\section*{Appendix: Short reproduction notes}
If you want exact environment pins, install with (CPU-only example):
\begin{verbatim}
python3 -m pip install --user numpy
python3 -m pip install --user torch torchvision torchaudio --index-url https://download.pytorch.org/whl/cpu
\end{verbatim}

\vspace{6pt}
\noindent{}If you prefer, I can add a short script that computes a confusion matrix and saves training/validation curves as PNGs to include in the PDF.

\end{document}
